\documentclass[12pt,a4paper]{report}

% ── Packages ──────────────────────────────────────────────────────────────────
\usepackage[T1]{fontenc}
\usepackage[utf8]{inputenc}
\usepackage[margin=2.5cm]{geometry}
\usepackage{lmodern}
\usepackage{microtype}
\usepackage{hyperref}
\usepackage{booktabs}
\usepackage{longtable}
\usepackage{array}
\usepackage{enumitem}
\usepackage{titlesec}
\usepackage{fancyhdr}
\usepackage{xcolor}
\usepackage{graphicx}
\usepackage{tcolorbox}
\usepackage{parskip}

% ── Colours ───────────────────────────────────────────────────────────────────
\definecolor{broadblue}{RGB}{0,60,120}
\definecolor{lightgray}{RGB}{245,245,245}
\definecolor{warnyellow}{RGB}{255,243,205}
\definecolor{warnborder}{RGB}{204,153,0}
\definecolor{noteblue}{RGB}{219,234,254}
\definecolor{noteborder}{RGB}{59,130,246}

% ── Hyperref ──────────────────────────────────────────────────────────────────
\hypersetup{
  colorlinks=true,
  linkcolor=broadblue,
  urlcolor=broadblue,
  pdftitle={DashyDashboard Deployment Guide},
  pdfauthor={Broadridge},
  pdfsubject={IIS Deployment on Windows Server}
}

% ── Headers / footers ─────────────────────────────────────────────────────────
\pagestyle{fancy}
\fancyhf{}
\fancyhead[L]{\small\color{broadblue}DashyDashboard — Deployment Guide}
\fancyhead[R]{\small\color{broadblue}Broadridge}
\fancyfoot[C]{\small\thepage}
\renewcommand{\headrulewidth}{0.4pt}

% ── Section formatting ────────────────────────────────────────────────────────
\titleformat{\chapter}[display]
  {\normalfont\Large\bfseries\color{broadblue}}
  {Chapter \thechapter}{10pt}{\Huge}
\titlespacing*{\chapter}{0pt}{20pt}{20pt}

\titleformat{\section}
  {\normalfont\large\bfseries\color{broadblue}}
  {\thesection}{1em}{}

% ── Custom boxes ──────────────────────────────────────────────────────────────
\tcbuselibrary{skins,breakable}

\newtcolorbox{notebox}{
  colback=noteblue,
  colframe=noteborder,
  fonttitle=\bfseries,
  title=Note,
  breakable
}

\newtcolorbox{warnbox}{
  colback=warnyellow,
  colframe=warnborder,
  fonttitle=\bfseries,
  title=Important,
  breakable
}

% ── Misc ──────────────────────────────────────────────────────────────────────
\setlist[enumerate]{itemsep=4pt,parsep=0pt}
\setlist[itemize]{itemsep=4pt,parsep=0pt}

% ═════════════════════════════════════════════════════════════════════════════
\begin{document}

% ── Title page ────────────────────────────────────────────────────────────────
\begin{titlepage}
  \centering
  \vspace*{3cm}
  {\color{broadblue}\rule{\linewidth}{1.5pt}}\\[0.6cm]
  {\Huge\bfseries\color{broadblue} DashyDashboard}\\[0.4cm]
  {\LARGE Deployment Guide}\\[0.3cm]
  {\color{broadblue}\rule{\linewidth}{1.5pt}}\\[1cm]
  {\large ASP.NET 7 + React on IIS / SQL Server}\\[2cm]
  {\large\itshape Broadridge Financial Solutions}\\[0.5cm]
  {\large June 2026}\\[3cm]
  \vfill
  {\small
    \begin{tabular}{ll}
      \textbf{Document type:}   & Deployment / Operations Guide \\
      \textbf{Audience:}        & Windows Server Administrators \\
      \textbf{Classification:}  & Internal \\
    \end{tabular}
  }
\end{titlepage}

% ── Table of contents ─────────────────────────────────────────────────────────
\tableofcontents
\clearpage

% ═════════════════════════════════════════════════════════════════════════════
\chapter{Prerequisites}

Before beginning the deployment, ensure the following are in place on the target server.

\section{Operating System}
\begin{itemize}
  \item Windows Server 2019 or Windows Server 2022 (Datacenter or Standard edition)
  \item Windows 10 or Windows 11 Pro is also supported for development/test deployments
\end{itemize}

\section{SQL Server}
\begin{itemize}
  \item Any edition of Microsoft SQL Server must be installed and running (SQL Server 2019 or later recommended; SQL Server Express is acceptable)
  \item SQL Server Management Studio (SSMS) must be available to run the database setup script
  \item You must be able to connect to SQL Server with \texttt{Windows Authentication} from the target machine
\end{itemize}

\section{Account Permissions}
\begin{itemize}
  \item You must be logged in as, or be able to run programs as, a \textbf{Local Administrator}
  \item The account running SSMS must have \textbf{sysadmin} rights on the SQL Server instance (or at least the ability to create databases and logins)
\end{itemize}

\section{Deployment Package}
The deployment package (the folder containing this document) must include:
\begin{itemize}
  \item \texttt{app\textbackslash} --- the published ASP.NET application files (including \texttt{wwwroot\textbackslash} with the React front-end)
  \item \texttt{setup.sql} --- the database creation script
  \item This document
\end{itemize}

\begin{notebox}
No internet connectivity is required on the target server after the .NET 7 Hosting Bundle has been installed. All application binaries are self-contained in the \texttt{app\textbackslash} folder.
\end{notebox}

% ═════════════════════════════════════════════════════════════════════════════
\chapter{Install IIS (Web Server Role)}

Internet Information Services (IIS) is the web server that will host the application. If IIS is already installed on this server, skip to Chapter~\ref{ch:hostingbundle}.

\begin{enumerate}
  \item Open \textbf{Server Manager}. It usually opens automatically after login; if not, click the Windows Start button and type \emph{Server Manager}.

  \item In the Server Manager dashboard, click \textbf{Manage} (top-right menu) and then select \textbf{Add Roles and Features}.

  \item The \emph{Add Roles and Features Wizard} opens. On the \emph{Before You Begin} page, click \textbf{Next}.

  \item On the \emph{Installation Type} page, select \textbf{Role-based or Feature-based installation} and click \textbf{Next}.

  \item On the \emph{Server Selection} page, your local server should already be highlighted. Click \textbf{Next}.

  \item On the \emph{Server Roles} page, locate \textbf{Web Server (IIS)} in the list and tick its checkbox.
  \begin{itemize}
    \item A dialog box will appear asking to add required features. Click \textbf{Add Features}.
    \item Click \textbf{Next}.
  \end{itemize}

  \item On the \emph{Features} page, leave all defaults as-is and click \textbf{Next}.

  \item On the \emph{Web Server Role (IIS)} information page, click \textbf{Next}.

  \item On the \emph{Role Services} page, expand \textbf{Security} and ensure the following are ticked:
  \begin{itemize}
    \item \textbf{Windows Authentication}
    \item Request Filtering (ticked by default)
  \end{itemize}
  Click \textbf{Next}.

  \item On the \emph{Confirmation} page, click \textbf{Install}.

  \item Wait for the installation progress bar to complete. When it reads \emph{Installation succeeded}, click \textbf{Close}.
\end{enumerate}

\begin{notebox}
After installation, you can verify IIS is running by opening a browser on the server and navigating to \texttt{http://localhost}. You should see the default IIS welcome page (blue IIS logo).
\end{notebox}

% ═════════════════════════════════════════════════════════════════════════════
\chapter{Install the .NET 7 Hosting Bundle}
\label{ch:hostingbundle}

The .NET 7 Hosting Bundle installs the ASP.NET Core runtime and the IIS integration module (AspNetCoreModuleV2) that allows IIS to host .NET applications.

\section{Download}
\begin{enumerate}
  \item On a machine with internet access, open a browser and search for: \emph{dotnet 7 hosting bundle download}
  \item Navigate to the official Microsoft .NET download page (microsoft.com/net/download or dotnet.microsoft.com).
  \item Under the \textbf{.NET 7} section, locate the \textbf{Hosting Bundle} for Windows and download the installer (\texttt{dotnet-hosting-7.x.x-win.exe}).
  \item Transfer the installer file to the target server if you downloaded it on a different machine.
\end{enumerate}

\section{Install}
\begin{enumerate}
  \item Right-click the installer file and select \textbf{Run as administrator}.
  \item Accept the license agreement and click \textbf{Install}.
  \item Leave all options at their defaults and allow the installer to complete.
  \item Click \textbf{Close} when the installation finishes.
\end{enumerate}

\section{Register with IIS}
\begin{warnbox}
If IIS was installed \emph{before} the Hosting Bundle, you must restart IIS to register the
ASP.NET Core Module. Skip this step if you installed the Hosting Bundle before IIS.
\end{warnbox}

\begin{enumerate}
  \item Click Start, type \textbf{IIS Manager}, and open \emph{Internet Information Services
        (IIS) Manager}.
  \item In the left pane, click on the server node (the topmost item with a computer icon).
  \item In the centre pane, double-click \textbf{Modules}.
  \item Scroll through the list and confirm an entry named \textbf{AspNetCoreModuleV2} exists.
  \item If it is \emph{not} listed, the ASP.NET Core Module must be registered. Choose
        whichever of the following options is available on your server:\nopagebreak
  \begin{enumerate}[label=\alph*.]
    \item \textbf{Recycle the application pool (recommended on production servers)} ---
          in IIS Manager, expand \textbf{Application Pools}, right-click
          \texttt{DashyDashboardPool} and select \textbf{Recycle}. Then refresh the
          Modules list.
    \item \textbf{Full IIS reset} (only if you have the necessary permissions) ---
          open a Command Prompt as Administrator, type \texttt{iisreset} and press Enter.
          Wait for \emph{Internet services successfully restarted}, then return to IIS
          Manager, refresh the Modules list, and confirm \textbf{AspNetCoreModuleV2} appears.
  \end{enumerate}
\end{enumerate}

% ═════════════════════════════════════════════════════════════════════════════
\chapter{Create the Database}

\begin{enumerate}
  \item Open \textbf{SQL Server Management Studio (SSMS)} from the Start menu.

  \item In the \emph{Connect to Server} dialog:
  \begin{itemize}
    \item Server type: \textbf{Database Engine}
    \item Server name: enter the name of your SQL Server instance (e.g., \texttt{(local)} or \texttt{.\textbackslash SQLEXPRESS} for Express)
    \item Authentication: \textbf{Windows Authentication}
    \item Click \textbf{Connect}
  \end{itemize}

  \item In the \emph{Object Explorer} panel (left side), click \textbf{File} in the menu bar and select \textbf{Open → File\ldots} (or press \texttt{Ctrl+O}).

  \item Browse to the deployment package folder and select \textbf{setup.sql}. Click \textbf{Open}.

  \item The SQL script will open in a query window. Click \textbf{Execute} in the toolbar (or press \textbf{F5}).

  \item Wait for the script to finish. The \emph{Messages} tab at the bottom should show output ending with lines similar to:
  \begin{quote}
    \texttt{(N row(s) affected)}\\
    \texttt{Command(s) completed successfully.}
  \end{quote}

  \item In \emph{Object Explorer}, right-click \textbf{Databases} and select \textbf{Refresh}. Confirm \textbf{ClientsAppAttestation} now appears in the list.
\end{enumerate}

\begin{notebox}
If the script reports that the database already exists, it is safe to re-run as long as the script uses \texttt{IF NOT EXISTS} guards. If you need a clean install, drop the existing database first via right-click → Delete in Object Explorer.
\end{notebox}

% ═════════════════════════════════════════════════════════════════════════════
\chapter{Grant SQL Access to the IIS App Pool}

The application runs under the IIS application pool identity (\texttt{IIS APPPOOL\textbackslash DashyDashboardPool}). This Windows account must have permission to read and write the database.

\begin{warnbox}
You must complete Chapter~\ref{ch:site} (creating the application pool named \texttt{DashyDashboardPool}) before this chapter, because the login name is derived from the pool name. If you prefer, create the application pool first (Chapter~\ref{ch:site}, steps 1--4) and then return here.
\end{warnbox}

\begin{enumerate}
  \item In SSMS \emph{Object Explorer}, expand your server node, then expand \textbf{Security}.

  \item Right-click \textbf{Logins} and select \textbf{New Login\ldots}

  \item On the \emph{General} page:
  \begin{itemize}
    \item In the \emph{Login name} field, type exactly: \texttt{IIS APPPOOL\textbackslash DashyDashboardPool}
    \item \textbf{Do not} use the Browse button --- type the name manually
    \item Leave \emph{Windows authentication} selected
  \end{itemize}

  \item On the left panel, click \textbf{User Mapping}.

  \item In the upper list of databases, tick the checkbox next to \textbf{ClientsAppAttestation}.

  \item In the lower panel (\emph{Database role membership for: ClientsAppAttestation}), tick the following roles:
  \begin{itemize}
    \item \textbf{db\_datareader}
    \item \textbf{db\_datawriter}
    \item \textbf{db\_ddladmin}
  \end{itemize}

  \item Click \textbf{OK} to create the login and user mapping.
\end{enumerate}

\begin{notebox}
\texttt{db\_ddladmin} is required on first launch because Entity Framework will apply database migrations automatically if the schema is not yet present or is out of date.
\end{notebox}

% ═════════════════════════════════════════════════════════════════════════════
\chapter{Deploy the Application Files}

\begin{enumerate}
  \item Open \textbf{Windows Explorer}.

  \item Navigate to \texttt{C:\textbackslash inetpub\textbackslash wwwroot} (this folder is created automatically by IIS).

  \item Right-click in an empty area and select \textbf{New → Folder}. Name it \textbf{DashyDashboard}.
  \begin{itemize}
    \item You may use a different path if your organisation's convention requires it (e.g., \texttt{D:\textbackslash Apps\textbackslash DashyDashboard}). Note the full path for use in Chapter~\ref{ch:site}.
  \end{itemize}

  \item Open the deployment package's \texttt{app\textbackslash} folder.

  \item Select all contents of the \texttt{app\textbackslash} folder (\textbf{Ctrl+A}) and copy them (\textbf{Ctrl+C}).

  \item Navigate back to \texttt{C:\textbackslash inetpub\textbackslash wwwroot\textbackslash DashyDashboard} and paste (\textbf{Ctrl+V}).

  \item (Optional but recommended) Inside the \texttt{DashyDashboard} folder, create a subfolder named \textbf{logs}. This is where IIS will write stdout logs if you ever enable them for troubleshooting.
\end{enumerate}

After this step, the folder structure should look like:
\begin{quote}
  \texttt{C:\textbackslash inetpub\textbackslash wwwroot\textbackslash DashyDashboard\textbackslash}\\
  \quad \texttt{DashyDashboard.Api.dll}\\
  \quad \texttt{web.config}\\
  \quad \texttt{appsettings.Production.json}\\
  \quad \texttt{wwwroot\textbackslash index.html}\\
  \quad \texttt{wwwroot\textbackslash assets\textbackslash \ldots}\\
  \quad \texttt{logs\textbackslash}
\end{quote}

% ═════════════════════════════════════════════════════════════════════════════
\chapter{Create the IIS Site and Application Pool}
\label{ch:site}

\section{Create the Application Pool}
\begin{enumerate}
  \item Open \textbf{IIS Manager} (Start → search \emph{IIS Manager}).

  \item In the left \emph{Connections} pane, expand your server name.

  \item Right-click \textbf{Application Pools} and select \textbf{Add Application Pool\ldots}

  \item In the dialog:
  \begin{itemize}
    \item \textbf{Name:} \texttt{DashyDashboardPool}
    \item \textbf{.NET CLR version:} \textbf{No Managed Code}
    \item \textbf{Managed pipeline mode:} \textbf{Integrated}
  \end{itemize}
  Click \textbf{OK}.
\end{enumerate}

\begin{notebox}
Select \emph{No Managed Code} because ASP.NET Core runs its own managed runtime. Selecting a .NET CLR version here would have no effect and could cause warnings.
\end{notebox}

\section{Create the Website}
\begin{enumerate}
  \item In the left pane, right-click \textbf{Sites} and select \textbf{Add Website\ldots}

  \item Fill in the fields:
  \begin{itemize}
    \item \textbf{Site name:} \texttt{DashyDashboard}
    \item \textbf{Application pool:} click \textbf{Select\ldots} and choose \texttt{DashyDashboardPool}
    \item \textbf{Physical path:} click the \textbf{\ldots} button and browse to the folder you created in Chapter~6 (e.g., \texttt{C:\textbackslash inetpub\textbackslash wwwroot\textbackslash DashyDashboard})
    \item \textbf{Port:} \texttt{8080} (or another available port if 8080 is already in use)
    \item Leave the IP address as \textbf{All Unassigned}
  \end{itemize}
  Click \textbf{OK}.

  \item IIS Manager will ask whether to stop the existing Default Web Site (if it is also using port 80). If you are using port 8080 you can leave the Default Web Site running.
\end{enumerate}

% ═════════════════════════════════════════════════════════════════════════════
\chapter{Enable Windows Authentication}

\begin{enumerate}
  \item In IIS Manager, click the \textbf{DashyDashboard} site in the left pane.

  \item In the centre pane (Features View), locate the \textbf{Authentication} icon (under the \emph{IIS} section) and double-click it.

  \item In the \emph{Authentication} list:
  \begin{itemize}
    \item Click \textbf{Windows Authentication}. In the right \emph{Actions} panel, click \textbf{Enable}.
    \item Click \textbf{Anonymous Authentication}. In the right \emph{Actions} panel, click \textbf{Disable}.
  \end{itemize}

  \item The final state should show:
  \begin{itemize}
    \item Anonymous Authentication: \textbf{Disabled}
    \item Windows Authentication: \textbf{Enabled}
  \end{itemize}
\end{enumerate}

\begin{warnbox}
If \emph{Windows Authentication} does not appear in the list, the Windows Authentication role service was not installed in Chapter~2. Return to Server Manager, modify the Web Server (IIS) role, and add \textbf{Security → Windows Authentication}.
\end{warnbox}

% ═════════════════════════════════════════════════════════════════════════════
\chapter{Verify the Deployment}

\begin{enumerate}
  \item On the server, open a web browser (Edge, Chrome, or Firefox).

  \item Navigate to: \texttt{http://localhost:8080}

  \item Your browser may display a Windows credentials prompt. Enter your Windows username and password.

  \item The \textbf{DashyDashboard} login page should appear in the browser.

  \item The application uses Windows Authentication. Once the browser has authenticated you, it passes your identity to the application, which resolves your account from the \texttt{Users} table in the database. If your Windows account has a corresponding entry, you will be logged in automatically.
\end{enumerate}

\begin{notebox}
When accessing the application from other machines on the network, use the server's hostname or IP address instead of \texttt{localhost}, for example: \texttt{http://servername:8080}. Ensure that port 8080 is open in the Windows Firewall (Windows Defender Firewall → Advanced Settings → Inbound Rules → New Rule → Port 8080 TCP).
\end{notebox}

% ═════════════════════════════════════════════════════════════════════════════
\chapter{Troubleshooting}

The table below lists the most common deployment problems and their solutions.

\renewcommand{\arraystretch}{1.4}
\begin{longtable}{>{\raggedright\arraybackslash}p{4.5cm} >{\raggedright\arraybackslash}p{4.5cm} >{\raggedright\arraybackslash}p{5.5cm}}
\toprule
\textbf{Symptom} & \textbf{Likely Cause} & \textbf{Fix} \\
\midrule
\endhead

HTTP 500.21 --- ANCM handler not found &
IIS was installed \emph{after} the Hosting Bundle &
Open a Command Prompt as Administrator and run \texttt{iisreset}. This registers AspNetCoreModuleV2 with IIS. \\

\midrule

HTTP 500.30 --- App failed to start &
Database connection failing, or a startup exception in the application &
Verify SQL Server is running. Open \texttt{appsettings.Production.json} in the deployment folder and confirm the \texttt{Default} connection string points to the correct server. Enable stdout logging temporarily: open \texttt{web.config}, set \texttt{stdoutLogEnabled="true"}, restart the site, and check the \texttt{logs\textbackslash} folder. \\

\midrule

HTTP 400 Bad Request &
App pool identity not granted access to SQL Server &
Complete Chapter~5 in full. Ensure the login name is typed as \texttt{IIS APPPOOL\textbackslash DashyDashboardPool} exactly. \\

\midrule

HTTP 401 Unauthorized &
Windows Authentication not enabled or Anonymous Authentication still enabled &
Complete Chapter~8. Verify the Authentication feature settings match the required state (Windows Auth Enabled, Anonymous Disabled). \\

\midrule

Page loads but shows only a blank white screen or ``Loading\ldots'' indefinitely &
React front-end files (\texttt{wwwroot\textbackslash}) not copied, or the API cannot reach the database &
Verify that \texttt{wwwroot\textbackslash index.html} exists in the deployment folder (Chapter~6). If it exists, open browser developer tools (F12) and check the Network tab for failing API requests. \\

\midrule

``Loading\ldots'' spinner never resolves after login &
Application pool stopped or crashed &
In IIS Manager, right-click \textbf{DashyDashboardPool} under Application Pools and select \textbf{Start}. If it immediately stops again, enable stdout logging and inspect the error. \\

\midrule

Windows credentials prompt loops endlessly &
Kerberos/NTLM negotiation issue (common when accessing via IP address rather than hostname) &
Access the application using the server's hostname instead of its IP address. Alternatively, add the site to Internet Explorer's \emph{Local Intranet} zone via Internet Options. \\

\midrule

Database exists but tables are missing &
\texttt{setup.sql} was not fully executed &
Re-open \texttt{setup.sql} in SSMS, select all text (\texttt{Ctrl+A}), and click Execute (F5). Check for any error messages in the \emph{Messages} tab. \\

\bottomrule
\end{longtable}

\section{Useful Locations}

\begin{itemize}
  \item \textbf{Application files:} \texttt{C:\textbackslash inetpub\textbackslash wwwroot\textbackslash DashyDashboard\textbackslash}
  \item \textbf{IIS Manager:} Start → search \emph{IIS Manager}
  \item \textbf{Windows Event Log:} Start → search \emph{Event Viewer} → Windows Logs → Application (ASP.NET Core errors are logged here)
  \item \textbf{IIS logs:} \texttt{C:\textbackslash inetpub\textbackslash logs\textbackslash LogFiles\textbackslash W3SVC\{n\}\textbackslash}
  \item \textbf{Connection string:} \texttt{app\textbackslash appsettings.Production.json}
\end{itemize}

% ═════════════════════════════════════════════════════════════════════════════
% PART 2 — Application Reference
% ═════════════════════════════════════════════════════════════════════════════
\part{Application Reference}

% ─────────────────────────────────────────────────────────────────────────────
\chapter{Application Overview}

DashyDashboard is a tool attestation management system for Broadridge. It records, tracks, and
reports whether associates actively used the client-facing tools to which they had access during
a given period.

\section{Purpose}

Each calendar month, an administrator opens a new \emph{attestation cycle}. During that cycle,
every associate whose account appears in the system receives a personalised list of the client
tools they were granted access to. The associate reviews each tool and confirms:

\begin{enumerate}
  \item Whether they actually had access to the tool during the cycle (the system pre-fills this
        based on the access records, but the associate may dispute it).
  \item Whether they used that tool during the cycle.
\end{enumerate}

Each attestation row can also carry a free-text remark. Once an associate has reviewed all their
tools they submit their attestation; the submission timestamp is recorded and the associate's
status updates to \emph{Submitted}.

\section{Role Summary}

\begin{itemize}
  \item \textbf{Associates} — log in, attest their own tools.
  \item \textbf{Managers} — see a progress dashboard for their direct reports and are alerted if
        any report disputes their access.
  \item \textbf{Administrators (Admin)} — full department-wide completion dashboard; can add
        new tools.
  \item \textbf{GFHs and GFH-Delegates} — department-level dashboard, same scope as Admin.
  \item \textbf{IFHs} — department-scoped reporting, configurable access level.
\end{itemize}

\section{Authentication}

The application relies entirely on \textbf{Windows Authentication} via IIS. No username or
password form is presented to the user. IIS authenticates the browser session using
Kerberos/NTLM and passes the Windows identity to the application, which resolves the
corresponding account from the \texttt{dbo.Users} table by matching the \texttt{WindowsID}
column. Every login attempt (successful or failed) is recorded in \texttt{dbo.LoginLogs}.

% ─────────────────────────────────────────────────────────────────────────────
\chapter{Database Schema}

The application database is named \textbf{ClientsAppAttestation}. It contains the ten tables
described below. All string primary keys that may carry leading zeros (such as client
identifiers) are stored as \texttt{VARCHAR} to preserve those characters exactly.

% ── 1. Departments ────────────────────────────────────────────────────────────
\section{dbo.Departments}

Holds the list of business departments. Every user, client, tool, and access record is linked
to a department, enabling scoped reporting for managers and GFHs.

\begin{longtable}{>{\ttfamily}p{4.5cm} >{\raggedright\arraybackslash}p{9.5cm}}
\toprule
\normalfont\textbf{Column} & \textbf{Notes} \\
\midrule
\endhead
DepartmentID & \texttt{INT IDENTITY} — primary key, auto-incremented. \\
DepartmentName & \texttt{VARCHAR(150) NOT NULL} — human-readable name of the department. \\
\bottomrule
\end{longtable}

% ── 2. Users ──────────────────────────────────────────────────────────────────
\section{dbo.Users}

Stores one row per associate. \texttt{AssociateID} is the business primary key and is treated as
a string throughout the system so that alphanumeric and zero-padded identifiers are preserved
without corruption.

\begin{longtable}{>{\ttfamily}p{4.5cm} >{\raggedright\arraybackslash}p{9.5cm}}
\toprule
\normalfont\textbf{Column} & \textbf{Notes} \\
\midrule
\endhead
AssociateID & \texttt{VARCHAR(50) NOT NULL} — primary key (not an identity column). \\
FirstName   & \texttt{VARCHAR(100)} \\
LastName    & \texttt{VARCHAR(100)} \\
EMailAddr   & \texttt{VARCHAR(255)} — corporate email address. \\
WindowsID   & \texttt{VARCHAR(200)} — Windows domain account used to match the IIS
              authenticated identity at login. \\
ManagerID   & \texttt{VARCHAR(50) NULL} — self-referencing foreign key to
              \texttt{dbo.Users.AssociateID}. Any user whose \texttt{AssociateID} appears as
              another user's \texttt{ManagerID} is automatically treated as a manager by the
              application. \\
DepartmentID & \texttt{INT} — foreign key to \texttt{dbo.Departments}. \\
IsActive    & \texttt{BIT NOT NULL DEFAULT 1} — set to \texttt{0} to deactivate an account
              without deleting it. \\
\bottomrule
\end{longtable}

% ── 3. Clients ────────────────────────────────────────────────────────────────
\section{dbo.Clients}

Holds the client entities whose tools associates must attest. \texttt{ClientID} is stored as
\texttt{VARCHAR} so that identifiers containing leading zeros or non-numeric characters (for
example \texttt{BARCLAYS} or \texttt{DTC-US}) are preserved as-is.

\begin{longtable}{>{\ttfamily}p{4.5cm} >{\raggedright\arraybackslash}p{9.5cm}}
\toprule
\normalfont\textbf{Column} & \textbf{Notes} \\
\midrule
\endhead
ClientID     & \texttt{VARCHAR(50) NOT NULL} — primary key (not an identity column). \\
ClientName   & \texttt{VARCHAR(255) NOT NULL} \\
DepartmentID & \texttt{INT} — foreign key to \texttt{dbo.Departments}. \\
IsActive     & \texttt{BIT NOT NULL DEFAULT 1} \\
\bottomrule
\end{longtable}

% ── 4. ClientTools ────────────────────────────────────────────────────────────
\section{dbo.ClientTools}

Defines the set of tools that belong to each client. A single tool is uniquely identified by
the combination of its client and a numeric tool identifier.

\begin{longtable}{>{\ttfamily}p{4.5cm} >{\raggedright\arraybackslash}p{9.5cm}}
\toprule
\normalfont\textbf{Column} & \textbf{Notes} \\
\midrule
\endhead
ClientID     & \texttt{VARCHAR(50) NOT NULL} — part of composite primary key;
               foreign key to \texttt{dbo.Clients}. \\
ToolID       & \texttt{INT NOT NULL} — part of composite primary key. \\
ToolName     & \texttt{NVARCHAR(200) NOT NULL} — display name of the tool. \\
DepartmentID & \texttt{INT} — foreign key to \texttt{dbo.Departments}. \\
\bottomrule
\end{longtable}

\textbf{Primary key:} composite \texttt{(ClientID, ToolID)}.

% ── 5. UsersToolAccess ────────────────────────────────────────────────────────
\section{dbo.UsersToolAccess}

Records which associate has access to which client tool, when that access was granted, and
whether it is a primary or secondary access tier. One row represents one
associate-client-tool access grant.

\begin{longtable}{>{\ttfamily}p{4.5cm} >{\raggedright\arraybackslash}p{9.5cm}}
\toprule
\normalfont\textbf{Column} & \textbf{Notes} \\
\midrule
\endhead
AssociateID  & \texttt{VARCHAR(50) NOT NULL} — part of composite primary key;
               foreign key to \texttt{dbo.Users}. \\
ClientID     & \texttt{VARCHAR(50) NOT NULL} — part of composite primary key;
               foreign key to \texttt{dbo.Clients}. \\
ToolID       & \texttt{INT NOT NULL} — part of composite primary key. \\
Tier         & \texttt{NVARCHAR(50) NOT NULL} — \texttt{Primary} or \texttt{Secondary}. \\
GivenDate    & \texttt{DATE NOT NULL} — date access was granted. \\
AccessTo     & \texttt{DATE NULL} — date access is scheduled to expire; \texttt{NULL} means
               open-ended. \\
GrantedBy    & \texttt{VARCHAR(50) NULL} — foreign key to \texttt{dbo.Users.AssociateID};
               the associate who granted this access. \\
DepartmentID & \texttt{INT} — foreign key to \texttt{dbo.Departments}. \\
\bottomrule
\end{longtable}

\textbf{Primary key:} composite \texttt{(AssociateID, ClientID, ToolID)}.

% ── 6. Cycles ─────────────────────────────────────────────────────────────────
\section{dbo.Cycles}

Defines each attestation cycle. The application reads the most recent active cycle
automatically; no application restart is required when a new cycle is added.

\begin{longtable}{>{\ttfamily}p{4.5cm} >{\raggedright\arraybackslash}p{9.5cm}}
\toprule
\normalfont\textbf{Column} & \textbf{Notes} \\
\midrule
\endhead
CycleID   & \texttt{INT IDENTITY} — primary key, auto-incremented. \\
CycleName & \texttt{NVARCHAR(100) NOT NULL} — descriptive label, e.g.\ ``Q2 2026 Attestation''. \\
StartDate & \texttt{DATE NOT NULL} — first day of the attestation window. \\
EndDate   & \texttt{DATE NOT NULL} — last day of the attestation window. \\
DueDate   & \texttt{DATE NOT NULL} — deadline by which all associates must have submitted. \\
\bottomrule
\end{longtable}

% ── 7. ToolCycleAttestation ───────────────────────────────────────────────────
\section{dbo.ToolCycleAttestation}

The central fact table. One row exists per associate-tool combination within a cycle. Rows are
created by the application when a cycle opens and are updated as associates attest each tool.

\begin{longtable}{>{\ttfamily}p{4.5cm} >{\raggedright\arraybackslash}p{9.5cm}}
\toprule
\normalfont\textbf{Column} & \textbf{Notes} \\
\midrule
\endhead
CycleID           & \texttt{INT NOT NULL} — part of composite primary key;
                    foreign key to \texttt{dbo.Cycles}. \\
AssociateID       & \texttt{VARCHAR(50) NOT NULL} — part of composite primary key;
                    foreign key to \texttt{dbo.Users}. \\
ClientID          & \texttt{VARCHAR(50) NOT NULL} — part of composite primary key;
                    foreign key to \texttt{dbo.Clients}. \\
ToolID            & \texttt{INT NOT NULL} — part of composite primary key. \\
AttestationStatus & \texttt{NVARCHAR(50) NOT NULL DEFAULT 'Pending'} — transitions to
                    \texttt{Submitted} when the associate submits. \\
UsedThisCycle     & \texttt{BIT NULL} — \texttt{1} = used, \texttt{0} = not used,
                    \texttt{NULL} = not yet answered. \\
HadAccess         & \texttt{BIT NOT NULL DEFAULT 1} — system-populated from
                    \texttt{dbo.UsersToolAccess}; the associate may set this to \texttt{0}
                    to dispute that they had access. \\
Remarks           & \texttt{NVARCHAR(500) NULL} — free-text comment from the associate. \\
SubmittedAt       & \texttt{DATETIME NULL} — populated when the associate submits. \\
\bottomrule
\end{longtable}

\textbf{Primary key:} composite \texttt{(CycleID, AssociateID, ClientID, ToolID)}.

% ── 8. AttestationLogs ────────────────────────────────────────────────────────
\section{dbo.AttestationLogs}

Provides a summary-level audit trail of each associate's submission event. One row is written
per associate per cycle at the moment of submission.

\begin{longtable}{>{\ttfamily}p{4.5cm} >{\raggedright\arraybackslash}p{9.5cm}}
\toprule
\normalfont\textbf{Column} & \textbf{Notes} \\
\midrule
\endhead
LogID       & \texttt{INT IDENTITY} — primary key, auto-incremented. \\
CycleID     & \texttt{INT NOT NULL} — foreign key to \texttt{dbo.Cycles}. \\
AssociateID & \texttt{VARCHAR(50) NOT NULL} — foreign key to \texttt{dbo.Users}. \\
SubmittedAt & \texttt{DATETIME NOT NULL} \\
ToolCount   & \texttt{INT NOT NULL} — number of tools included in the submission. \\
Summary     & \texttt{NVARCHAR(100) NULL} — short human-readable summary generated at
              submission time. \\
\bottomrule
\end{longtable}

% ── 9. SuperUsers ─────────────────────────────────────────────────────────────
\section{dbo.SuperUsers}

Grants elevated roles to specific associates. An associate may hold more than one role (for
example, both \texttt{GFH} and \texttt{IFH}) by having multiple rows, one per role. The unique
index on \texttt{(AssociateID, RoleName, DepartmentID)} prevents duplicate grants.

\begin{longtable}{>{\ttfamily}p{4.5cm} >{\raggedright\arraybackslash}p{9.5cm}}
\toprule
\normalfont\textbf{Column} & \textbf{Notes} \\
\midrule
\endhead
SuperUserID  & \texttt{INT IDENTITY} — primary key, auto-incremented. \\
AssociateID  & \texttt{VARCHAR(50) NOT NULL} — foreign key to \texttt{dbo.Users}. \\
RoleName     & \texttt{VARCHAR(50) NOT NULL} — one of \texttt{Admin}, \texttt{GFH},
               \texttt{GFHDelegate}, \texttt{IFH}. \\
DepartmentID & \texttt{INT NOT NULL} — foreign key to \texttt{dbo.Departments}; scopes the
               role to a department. \texttt{Admin} role ignores this constraint and sees all
               departments. \\
AccessLevel  & \texttt{VARCHAR(50) NOT NULL} — \texttt{Full}, \texttt{ReadOnly}, or
               \texttt{Approver}. \\
IsActive     & \texttt{BIT NOT NULL DEFAULT 1} — set to \texttt{0} to revoke the role without
               deleting the row. \\
\bottomrule
\end{longtable}

\textbf{Unique index:} \texttt{(AssociateID, RoleName, DepartmentID)}.

% ── 10. LoginLogs ─────────────────────────────────────────────────────────────
\section{dbo.LoginLogs}

Records every login attempt. \texttt{AssociateID} is nullable so that failed attempts (where
the Windows identity did not match any user in \texttt{dbo.Users}) are still captured with a
\texttt{Status} of \texttt{UserNotFound}.

\begin{longtable}{>{\ttfamily}p{4.5cm} >{\raggedright\arraybackslash}p{9.5cm}}
\toprule
\normalfont\textbf{Column} & \textbf{Notes} \\
\midrule
\endhead
LoginLogID  & \texttt{INT IDENTITY} — primary key, auto-incremented. \\
AssociateID & \texttt{VARCHAR(50) NULL} — foreign key to \texttt{dbo.Users}; \texttt{NULL}
              when no matching user was found. \\
LoginTime   & \texttt{DATETIME NOT NULL} — timestamp of the attempt. \\
Status      & \texttt{NVARCHAR(50) NOT NULL} — \texttt{Success} or \texttt{UserNotFound}. \\
\bottomrule
\end{longtable}

% ─────────────────────────────────────────────────────────────────────────────
\chapter{User Roles}

Access to application features is determined by the user's role. The base role (Associate) is
implicit for every account in \texttt{dbo.Users}. Elevated roles are stored in
\texttt{dbo.SuperUsers} with the exception of Manager, which is derived automatically from the
\texttt{ManagerID} hierarchy in \texttt{dbo.Users}.

\section{Associate}

The default role. Any account present in \texttt{dbo.Users} with \texttt{IsActive = 1} is an
Associate.

\begin{itemize}
  \item Logs in via Windows Authentication. The application resolves identity by matching
        the IIS-supplied Windows account name against \texttt{dbo.Users.WindowsID}.
  \item Can view and submit attestations for their own tool access in the current cycle.
  \item Attestation is a two-step process per tool:
  \begin{enumerate}
    \item \textbf{Had Access?} — the system pre-populates this as true based on
          \texttt{dbo.UsersToolAccess}. The associate may set it to false to dispute the
          access record.
    \item \textbf{Used This Cycle?} — this field is disabled if ``Had Access'' is false.
  \end{enumerate}
  \item The Remarks field is always editable regardless of the other field values.
\end{itemize}

\section{Manager}

No \texttt{SuperUsers} row is required. An associate becomes a Manager automatically when their
\texttt{AssociateID} appears as the \texttt{ManagerID} of at least one other active user.

\begin{itemize}
  \item Sees a team dashboard listing each direct report and their attestation completion
        status for the current cycle.
  \item Receives a mismatch notification if any direct report has disputed their access
        (\texttt{HadAccess = 0}); the manager can drill in to view the dispute details.
  \item A Manager is also an Associate and can attest their own tools in the same session.
\end{itemize}

\section{Admin}

Stored in \texttt{dbo.SuperUsers} with \texttt{RoleName = 'Admin'}.

\begin{itemize}
  \item Full department-wide overview dashboard: completion rates and per-client breakdowns
        across all departments.
  \item Can add new tools via the \textbf{Add Tool} button. The form requires selecting a
        \textbf{Client}, a \textbf{Department}, and entering a \textbf{Tool Name}. Admins
        may add tools to any client in any department.
  \item The \texttt{DepartmentID} column on the Admin's \texttt{SuperUsers} row is recorded
        but does not restrict visibility --- Admins see all departments.
\end{itemize}

\section{GFH (Global Functional Head)}

Stored in \texttt{dbo.SuperUsers} with \texttt{RoleName = 'GFH'}.

\begin{itemize}
  \item One GFH is displayed per department on the Admin dashboard.
  \item Has the same dashboard access as Admin: completion rates and per-client breakdowns.
  \item Can add new tools via the \textbf{Add Tool} button, scoped to their own department.
  \item Intended for the senior leader who is accountable for a department's attestation
        completion.
\end{itemize}

\section{GFH-Delegate}

Stored in \texttt{dbo.SuperUsers} with \texttt{RoleName = 'GFHDelegate'}.

\begin{itemize}
  \item Carries identical system access to GFH --- same dashboard view, same department scope,
        and the same ability to add tools via the \textbf{Add Tool} button (scoped to their
        department).
  \item Used when a GFH delegates day-to-day attestation oversight to a colleague. The
        delegate acts on behalf of the GFH but is not themselves the GFH.
  \item Multiple GFH-Delegates may be assigned to the same department by inserting one
        \texttt{dbo.SuperUsers} row per delegate.
  \item To assign a GFH-Delegate: insert a row into \texttt{dbo.SuperUsers} with
        \texttt{RoleName = 'GFHDelegate'}, the delegate's \texttt{AssociateID}, and the
        appropriate \texttt{DepartmentID}.
\end{itemize}

\section{IFH (Internal Functional Head)}

Stored in \texttt{dbo.SuperUsers} with \texttt{RoleName = 'IFH'}.

\begin{itemize}
  \item Currently reserved for department-scoped reporting.
  \item The exact access level is configurable via the \texttt{AccessLevel} field:
        \texttt{Full}, \texttt{ReadOnly}, or \texttt{Approver}.
\end{itemize}

% ─────────────────────────────────────────────────────────────────────────────
\chapter{Adding a User to a Role}

Role assignments are managed directly in SQL Server Management Studio. No application
deployment or restart is required.

\section{Prerequisites}

\begin{itemize}
  \item The associate must already have a row in \texttt{dbo.Users} with \texttt{IsActive = 1}.
  \item You must know the correct \texttt{DepartmentID} integer from \texttt{dbo.Departments}.
  \item You must be connected to the \textbf{ClientsAppAttestation} database in SSMS with
        sufficient permissions to edit table rows (\texttt{db\_datawriter} or higher).
\end{itemize}

\section{Steps}

\begin{enumerate}
  \item Open \textbf{SSMS} and connect to the database server using Windows Authentication.
  \item In \emph{Object Explorer}, expand \textbf{Databases} $\to$
        \textbf{ClientsAppAttestation} $\to$ \textbf{Tables}.
  \item Right-click \textbf{dbo.SuperUsers} and select \textbf{Edit Top 200 Rows}.
  \item Scroll to the first empty row at the bottom of the grid.
  \item Enter values for the following columns:
  \begin{itemize}
    \item \textbf{AssociateID} — the associate's ID exactly as it appears in
          \texttt{dbo.Users}. Must not be null.
    \item \textbf{RoleName} — one of: \texttt{Admin}, \texttt{GFH}, \texttt{GFHDelegate},
          \texttt{IFH}.
    \item \textbf{DepartmentID} — the integer department ID from \texttt{dbo.Departments}.
    \item \textbf{AccessLevel} — \texttt{Full} for Admin/GFH/GFHDelegate; \texttt{Full},
          \texttt{ReadOnly}, or \texttt{Approver} for IFH.
    \item \textbf{IsActive} — \texttt{1}.
  \end{itemize}
  \item Press \textbf{Tab} or click another row to commit the new entry.
\end{enumerate}

\begin{notebox}
The change takes effect immediately. The associate will see the elevated role on their next
login without any application restart.
\end{notebox}

\section{Removing a Role}

Set \texttt{IsActive = 0} on the relevant row rather than deleting it. This deactivates the
role while preserving the audit trail of when and to whom it was previously granted.

% ─────────────────────────────────────────────────────────────────────────────
\chapter{Attestation Cycle Management}

\section{Opening a New Cycle}

Cycles are managed directly in the database via SSMS. The application reads the most recently
opened cycle automatically; no code change or restart is required.

\begin{enumerate}
  \item Open \textbf{SSMS} and connect to the database server.
  \item In \emph{Object Explorer}, expand \textbf{Databases} $\to$
        \textbf{ClientsAppAttestation} $\to$ \textbf{Tables}.
  \item Right-click \textbf{dbo.Cycles} and select \textbf{Edit Top 200 Rows}.
  \item Scroll to the first empty row and enter:
  \begin{itemize}
    \item \textbf{CycleName} — a descriptive label, for example
          \texttt{Q2 2026 Attestation}.
    \item \textbf{StartDate} — first day of the attestation window (format
          \texttt{YYYY-MM-DD}).
    \item \textbf{EndDate} — last day of the attestation window.
    \item \textbf{DueDate} — deadline by which all associates must submit.
  \end{itemize}
  \item Press \textbf{Tab} or click another row to commit.
\end{enumerate}

\begin{notebox}
Associates see the new cycle and their updated tool list on their next login. No data from
previous cycles is overwritten; each cycle's attestation records are fully independent.
\end{notebox}

\section{Cycle Lifecycle}

When a cycle is active, associates can:

\begin{itemize}
  \item View the full list of client tools they have access to for that cycle.
  \item Set ``Had Access'' and ``Used This Cycle'' for each tool.
  \item Add or edit remarks at any time before the due date.
  \item Submit their attestation (sets \texttt{AttestationStatus = 'Submitted'} and records
        \texttt{SubmittedAt}).
\end{itemize}

A submission entry is also written to \texttt{dbo.AttestationLogs} recording the total number
of tools attested and a brief summary. Managers and Admins can monitor completion progress
in real time via their dashboards.

% ═════════════════════════════════════════════════════════════════════════════
\end{document}
